% ==============================================================================
% CHƯƠNG 4: QUY TRÌNH THỰC NGHIỆM 4 BƯỚC (PIPELINE 4 NOTEBOOKS)
% File: 05-chuong-4-thuc-nghiem.tex
% ==============================================================================

\chapter{Quy trình Thực nghiệm 4 Bước (Pipeline 4 Notebooks)}

\section{Notebook 01 — EDA \& Sơ đồ Phân chia Dữ liệu Stratified}

Notebook 01 bắt đầu bằng nguyên tắc bảo mật thông tin nghiêm ngặt: **Chia dữ liệu TRƯỚC KHI thực hiện bất kỳ phân tích hay trực quan hóa nào**, nhằm triệt tiêu hoàn toàn rủi ro rò rỉ dữ liệu (Data Leakage).

Dữ liệu ban đầu gồm 2 file thô: \texttt{full\_trainset.csv} ($926.669$ dòng dữ liệu quan sát) và \texttt{full\_testset.csv} ($181.669$ dòng dữ liệu ngẫu nhiên RCT). Notebook 01 thiết kế sơ đồ phân chia phân tầng (Stratified Split) theo $4$ tổ hợp nhãn $(is\_treat \times label)$:

\begin{equation}
    \text{Strata}_i = \text{str}(W_i) + \text{\_} + \text{str}(Y_i) \quad \in \{0\_0, 0\_1, 1\_0, 1\_1\}
\end{equation}

\begin{figure}[h!]
\centering
\begin{mdframed}[linecolor=lazadablue,backgroundcolor=codebg,linewidth=1pt]
\begin{verbatim}
full_trainset (926k) ──► train      (80% ~ 741.335 dòng) ──► EDA & Train
                     └─► val        (20% ~ 185.334 dòng) ──► Tuning Optuna (NB03)

full_testset  (181k) ──► rct_select (50% ~  90.834 dòng) ──► Chọn Mô hình Quán quân
                     └─► rct_holdout(50% ~  90.835 dòng) ──► NIÊM PHONG TUYỆT ĐỐI!
\end{verbatim}
\end{mdframed}
\caption{Sơ đồ phân chia 4 tập dữ liệu tại Notebook 01}
\end{figure}

Ngay sau khi phân chia, tập \texttt{val} và \texttt{rct\_holdout} bị xóa khỏi RAM và niêm phong trong đĩa cứng. File nhãn \texttt{split\_assignment.parquet} được lưu lại để tất cả các notebook sau sử dụng chung một sơ đồ chia nhất quán $100\%$.

\section{Notebook 02 — Kỹ nghệ Đặc trưng \& Kiểm tra Overlap}

Notebook 02 xử lý dữ liệu sạch và tạo lập các tính năng mới:
\begin{enumerate}[leftmargin=*]
    \item \textbf{Làm sạch dữ liệu:} Loại bỏ các cột hằng số ($\text{Variance} = 0$), các cột trùng lặp thông tin (băm hash so sánh) và bỏ biến nhị phân tuyến tính thừa. Từ $83$ cột thô thu về $71$ cột gốc sạch.
    \item \textbf{Kỹ nghệ 3 nhóm đặc trưng mới (Feature Engineering):}
    \begin{*itemize}
        \item \textit{Nhóm 1 (Tỉ lệ tương quan):} Các tỉ lệ giữa cặp đặc trưng tương quan cao (ví dụ: \texttt{fe\_ratio\_f1\_f2 = f1 / (f2 + 1)}).
        \item \textit{Nhóm 2 (Đặc trưng tương tác):} Tích giữa các biến hành vi (\texttt{fe\_inter\_f9\_f26 = f9 * f26}).
        \item \textit{Nhóm 3 (Thống kê gộp theo dòng):} Đếm số cột khác 0 trong Top đặc trưng (\texttt{fe\_nonzero\_top10}) và tổng cờ nhị phân (\texttt{fe\_flag\_sum}).
    \end{*itemize}
    Tổng số đặc trưng nâng từ $71$ lên **$76$ đặc trưng sạch**.
    \item \textbf{Kiểm tra Giả định Overlap (Positivity Check):} Huấn luyện mô hình Propensity score $\hat{e}(X)$ bằng LightGBM Classifier, vẽ biểu đồ phân phối xác suất và xác nhận 2 nhóm $W=1$ và $W=0$ có vùng chồng lấn xác suất rộng khắp.
\end{enumerate}

\section{Notebook 03 — Tinh chỉnh Tham số Optuna \& Log MLflow}

Để tối ưu hóa siêu tham số cho 6 mô hình mà không tốn cả ngày tính toán, Notebook 03 thực hiện:
\begin{itemize}[leftmargin=*]
    \item \textbf{Sub-sampling tiết kiệm thời gian:} Rút mẫu phân tầng $200.000$ dòng để Tune (\texttt{X\_tune}) và $100.000$ dòng làm Validation (\texttt{X\_val}).
    \item \textbf{Hàm mục tiêu DR-AUUC (DR-Qini) chuẩn hóa:} Thay vì dùng DR-MSE bị chi phối bởi phương sai nhiễu, Optuna chuyển sang tối ưu trực tiếp khả năng xếp hạng **DR-AUUC**:
    \begin{equation}
        \text{DR-AUUC} = \frac{\text{Area}(\text{Mô hình}) - \text{Area}(\text{Ngẫu nhiên})}{\text{Area}(\text{Tiên tri}) - \text{Area}(\text{Ngẫu nhiên})}
    \end{equation}
    \item \textbf{Winsorizing \& Overlap Trimming:} Winsorize $\tilde{\Gamma}$ ở phân vị $0.5\% / 99.5\%$ để triệt tiêu outlier cực trị. Tự động trim vùng không chồng lấn theo quy tắc **Crump et al. (2009)** ($[0.02, 0.10]$).
    \item \textbf{Transportability Check:} Mang mô hình DR-Learner sang dự đoán trên \texttt{rct\_select}, thực hiện Z-test kiểm định độ chênh lệch ATE và tính chỉ số MDD (Minimum Detectable Difference).
    \item \textbf{MLflow Experiment Tracking:} Lưu vết $66$ trials vào database \texttt{mlflow.db} và xuất bộ tham số tối ưu ra file \texttt{artifacts/best_params.json}.
\end{itemize}

\section{Notebook 04 — Benchmark Evaluation \& Ablation Study}

Notebook 04 đóng vai trò là đấu trường chung kết và đóng gói sản phẩm:
\begin{enumerate}[leftmargin=*]
    \item \textbf{Kiểm chứng Metric (4 Sanity Checks):} Kiểm thử bộ hàm tính Qini, AUUC, ATE qua 4 phép thử ngược có \texttt{assert} để đảm bảo công thức toán học đúng $100\%$ trước khi chấm điểm.
    \item \textbf{Huấn luyện lại (Full Re-train):} Nạp \texttt{best_params.json}, gộp $100\%$ dữ liệu quan sát ($train + val = 926.669$ dòng) để huấn luyện lại 6 mô hình.
    \item \textbf{Đấu trên \texttt{rct\_select} \& Chọn Quán quân:} Chấm điểm Qini/AUUC kèm Bootstrap 95\% (1.000 lần). Mô hình **DR-Learner** xuất sắc giành vị trí Quán quân.
    \item \textbf{Feature Importance \& Ablation Study:} 
    Đo tầm quan trọng đặc trưng qua Gain ở tầng cuối DR-Learner. Quét lưới 9 mốc $k \in [7, 12, 17, 22, 26, 31, 36, 41, 76]$. Chốt mốc $k=41$ đặc trưng đầu vào (tương ứng **36 cột gốc trên Redis**) là mốc an toàn tuyệt đối khi phục vụ (sai lệch CATE = 0).
    \item \textbf{Đánh giá 1 lần duy nhất trên \texttt{rct\_holdout}:} Mở tập niêm phong đúng 1 lần, báo cáo Qini Score $\mathbf{0.0291}$ [$0.0041, 0.0553$], ATE $\mathbf{0.00376}$. Chạy Mô phỏng chính sách (Policy Simulation) và Phân tích nhóm Persuadables.
    \item \textbf{Xuất Artifacts bàn giao:} Đóng gói \texttt{model.pkl}, \texttt{metadata.json} và \texttt{feature_contract.json}.
\end{enumerate}
